\section{Water, Bread, and the Visible Neighbor}

Had the Incarnation been merely a brief thirty-three-year episode followed by a return to pure disembodied spirituality, Kant's suspicion might still have retained some plausibility: historical concreteness would look like a temporary bridge meant to be crossed once and left behind. Yet the Ascension of Christ did not inaugurate a religion of pure interiority detached from matter, place, and gesture. On the contrary, the Church continues the incarnate economy along two inseparable concrete axes: the visible sacraments instituted by Christ, and the visible neighbor whom Christ identifies with himself.

\subsection{From the Visible Redeemer to the Sacraments}

Preaching on the Ascension in \textit{Sermon 74}, St.\ Leo the Great explains why the withdrawal of Christ's visible bodily presence from the disciples' eyes did not leave the Church bereft of tangible encounter:

\begin{quote}
And so that which till then was visible of our Redeemer was changed into a sacramental presence, and that faith might be more excellent and stronger, sight gave way to doctrine, the authority of which was to be accepted by believing hearts enlightened with rays from above.\endnote{St.\ Leo the Great, \textit{Sermon 74} (\textit{On the Lord's Ascension}, II), \S\,2 (NPNF, Second Series, vol.\ 12; CCEL dated electronic text, 2026-09-28). Feltoe notes the underlying Latin phrase \latin{in sacramenta transivit} in his apparatus.}
\end{quote}

Leo's concise Latin phrase---\latin{in sacramenta transivit}, passed over into the sacraments---captures the continuity of Christian particularity across time. What the crowds in Galilee encountered through the touch of Christ's hand, the voice that spoke forgiveness, and the table he shared with sinners did not dissolve into a vague spiritual memory after the fortieth day. It became accessible to every century and every continent through concrete material signs united to the living word of Christ.

St.\ Augustine states two patristic principles that illuminate how material elements become instruments of grace. In \textit{The City of God}, analyzing the relation between outward ritual and interior devotion, he writes: ``A sacrifice, therefore, is the visible sacrament or sacred sign of an invisible sacrifice.''\endnote{Augustine, \textit{The City of God}, X.5.} God has no need of bulls, incense, water, or bread; he instituted visible signs for our sake, so that our inward turning to him might take embodied form and unite a whole people in common worship. Thus, Augustine continues in the next chapter, ``This is the sacrifice of Christians: we, being many, are one body in Christ. And this also is the sacrifice which the Church continually celebrates in the sacrament of the altar, known to the faithful, in which she teaches that she herself is offered in the offering she makes to God.''\endnote{Augustine, \textit{The City of God}, X.6.}

In his eightieth \textit{Tractate on the Gospel of John}, commenting on the Lord's words ``Now ye are clean through the word which I have spoken unto you,'' Augustine explains why ordinary water does not remain mere water at the font:

\begin{quote}
Take away the word, and the water is neither more nor less than water. The word is added to the element, and there results the Sacrament, as if itself also a kind of visible word.\endnote{St.\ Augustine, \textit{Tractates on the Gospel of John}, LXXX.3, trans.\ John Gibb and James Innes, in \textit{Nicene and Post-Nicene Fathers}, First Series, vol.\ 7 (CCEL dated electronic text, 2026-09-17; file header cites New York: Christian Literature Publishing Co., 1886).}
\end{quote}

A ``visible word'': just as the eternal Word took visible flesh in Mary's womb, so the spoken word of faith joins itself to sensible matter---water, chrism, bread, wine, the laying on of hands---to bring the saving power of Christ's Passion directly into the sensory fabric of human life.

\subsection{Why Embodied Knowers Need Sensible Signs}

Why would a spiritual God choose to convey spiritual grace through wet water, olive oil, and wheaten bread? Critics since antiquity have appealed even to Scripture itself---``God is a spirit: and they that adore him must adore him in spirit and in truth'' (John 4:24)---to argue that material rites belong to an immature stage of religion.

St.\ Thomas Aquinas answers that objection directly in his treatise on the sacraments (\textit{ST} III, qq.\,60--61), grounding sacramental particularity in human nature itself. ``Divine wisdom provides for each thing according to its mode... Now it is part of man's nature to acquire knowledge of the intelligible from the sensible.''\endnote{Aquinas, \textit{ST} III, q.\,60, a.\,4, corpus.} Angels are pure immaterial intellects; they need no water or oil to be enlightened. Human beings are neither disembodied minds trapped in flesh nor mere animals; we are unified composites of rational soul and living body whose every act of understanding begins in the senses.

Accordingly, Aquinas offers three interlocking reasons why sacraments are fitting and necessary for human salvation. First, human nature itself is led from corporeal and sensible realities to spiritual and intelligible ones. Second, because sin disordered human affection precisely by bending the heart toward material creatures, ``the healing remedy should be given to a man so as to reach the part affected by disease.'' Third, because bodily action is natural to human life, sacramental worship turns our physical energy away from superstition and vice toward ordered holiness. Aquinas summarizes the threefold fruit: ``through the institution of the sacraments man, consistently with his nature, is instructed through sensible things; he is humbled, through confessing that he is subject to corporeal things, seeing that he receives assistance through them: and he is even preserved from bodily hurt, by the healthy exercise of the sacraments.''\endnote{Aquinas, \textit{ST} III, q.\,61, a.\,1, corpus.}

The second fruit is humility. The proud intellect wants an unmediated, private ascent to the Absolute that owes nothing to poured water, spoken absolution, or a fellow sinner standing in the person of Christ. Submitting to receive invisible grace through humble created matter cures the exact pride that first revolted against creaturely dependence. As the \textit{Catechism of the Catholic Church} synthesizes this patristic and Thomistic doctrine, the sacramental liturgy is the work of Christ and his Body in which cosmic signs, human gestures, and covenantal realities are assumed by the Holy Spirit to communicate the grace of the Paschal mystery to embodied believers.\endnote{\textit{Catechism of the Catholic Church} (Latin \latin{editio typica}, 1997), \S\S\,1084--1085, 1145--1152 (doctrinal paraphrase).}

At the same time, Aquinas preserves the sovereign freedom of the universal Cause against any mechanical or magical narrowing of grace. Addressing whether sacraments or their human ministers confine God's action, he writes plainly: ``God did not bind His power to the sacraments, so as to be unable to bestow the sacramental effect without conferring the sacrament.''\endnote{Aquinas, \textit{ST} III, q.\,64, a.\,7, corpus.} The Church is bound to the sacraments Christ instituted; God himself is never imprisoned by them. Where invincible ignorance or sudden necessity prevents a person from reaching the font or the altar, the sovereign Author of grace remains free to save by ways known to him. Particular signs establish a reliable, visible covenantal home for embodied human beings; they never erect a wall around divine mercy.

\subsection{The Eucharist: Real Presence Without Local Confinement}

The sacramental defense of particularity reaches its apex in the Holy Eucharist. Here, Catholic faith makes a claim so concrete that many who gladly welcome religious symbolism draw back in disbelief: under the consecrated appearances of bread and wine, Jesus Christ is present not merely as a figure, a memorial, or a virtual influence, but in the true reality of his Body, Blood, Soul, and Divinity.

Aquinas begins his exposition in \textit{ST} III, q.\,75, a.\,1 by placing the mystery candidly within the obedience of faith: ``The presence of Christ's true body and blood in this sacrament cannot be detected by sense, nor understanding, but by faith alone, which rests upon Divine authority.'' Yet once received in faith, its fittingness illuminates the deepest meaning of friendship and love:

\begin{quote}
And because it is the special feature of friendship to live together with friends, as the Philosopher says (Ethic. ix), He promises us His bodily presence as a reward, saying (Matt. 24:28): ``Where the body is, there shall the eagles be gathered together.'' Yet meanwhile in our pilgrimage He does not deprive us of His bodily presence; but unites us with Himself in this sacrament through the truth of His body and blood... Hence this sacrament is the sign of supreme charity, and the uplifter of our hope, from such familiar union of Christ with us.\endnote{Aquinas, \textit{ST} III, q.\,75, a.\,1, corpus.}
\end{quote}

Love does not rest content with sending messages across a distance; friendship seeks presence. In the Eucharist, the metaphysics of transubstantiation unites total concreteness with universal availability, overcoming the spatial limitation that normally separates particulars from one another.

An ordinary material body occupies one circumscribed place by its dimensive quantity; it cannot be simultaneously whole in Rome, Tokyo, and a village chapel in the Andes, nor can it be whole inside every broken fragment of a loaf. Critics naturally object that Christ's risen body is a true human body in heaven and therefore cannot be truly upon thousands of altars at once (\textit{ST} III, q.\,75, a.\,1, obj.\,3). Aquinas replies by distinguishing \emph{local presence} (a body's dimensions commensurate with a surrounding place) from \emph{sacramental, substantial presence}. Because the conversion of bread and wine terminates directly in the \emph{substance} of Christ's body and blood---while his bodily dimensions are present by real concomitance rather than by filling the host's spatial outline---``Christ's body is not in this sacrament in the same way as a body is in a place, which by its dimensions is commensurate with the place; but in a special manner which is proper to this sacrament. Hence we say that Christ's body is upon many altars, not as in different places, but `sacramentally.'\,''\endnote{Aquinas, \textit{ST} III, q.\,75, a.\,1, ad 3.}

For the same reason, answering the parallel objection that an organic body's dimensive quantity requires its parts to occupy distinct parts of place (\textit{ST} III, q.\,76, a.\,3, obj.\,3), Aquinas shows that ``Christ's body is in this sacrament substantively, that is, in the way in which substance is under dimensions, but not after the manner of dimensions... And therefore it is manifest that the entire Christ is under every part of the species of the bread, even while the host remains entire.''\endnote{Aquinas, \textit{ST} III, q.\,76, a.\,3, corpus.} Here the ancient opposition between the particular and the universal collapses completely. Nothing could be more concrete than this single consecrated host laid upon a believer's tongue; yet nothing could be less provincially confined, for the whole Christ is given undivided to every communicant on every altar across the face of the earth, without leaving the Father's right hand or diminishing his gift to anyone else.

\subsection{The Visible Neighbor as the Moral Test of Particularity}

If the flight from particularity distorts metaphysics and worship, it works an even deadlier corruption in ethics. Few illusions flatter the human heart more easily than an ardent, tearful devotion to ``humanity in general'' paired with impatience, contempt, or indifference toward the actual, inconvenient neighbor sitting across the table or shivering outside the gate. Abstract humanity never interrupts our schedule, smells of poverty, repeats dull stories, or demands costly forgiveness. Concrete human beings do.

The First Epistle of St.\ John exposes that evasion plainly:

\begin{quote}
Let us therefore love God: because God first hath loved us. If any man say: I love God, and hateth his brother; he is a liar. For he that loveth not his brother whom he seeth, how can he love God whom he seeth not? And this commandment we have from God, that he who loveth God love also his brother (1 John 4:19--21).
\end{quote}

John's question reverses the Platonic order of ascent. The Platonist treats the visible singular as a distraction from the invisible God; the Apostle treats the visible brother as the concrete test of whether our professed love for the invisible God is real or fraudulent. Our Lord himself fixes the same standard for the Last Judgment in Matthew 25. He does not ask whether the nations entertained elevated sentiments about universal benevolence. He cites six stubbornly physical, local acts---food for the hungry, drink for the thirsty, shelter for the stranger, clothing for the naked, visits to the sick and the imprisoned---and binds himself personally to each sufferer: ``Amen I say to you, as long as you did it to one of these my least brethren, you did it to me'' (Matt 25:35--36, 40).

In \textit{Homily 50 on Matthew}, St.\ John Chrysostom ties the Eucharistic table directly to the shivering beggar outside the church door. Because the very same Lord who spoke ``This is my body'' over the altar also spoke ``Inasmuch as ye did it not to one of the least of these, ye did it not to me,'' Chrysostom refuses to let liturgical splendor excuse neglect of concrete charity---not forbidding fitting adornment of the sanctuary, yet demanding the order of reality:

\begin{quote}
For what is the profit, when His table indeed is full of golden cups, but He perishes with hunger? First fill Him, being an hungered, and then abundantly deck out His table also. Dost thou make Him a cup of gold, while thou givest Him not a cup of cold water?... Do not therefore while adorning His house overlook thy brother in distress, for he is more properly a temple than the other.\endnote{St.\ John Chrysostom, \textit{Homilies on the Gospel of St.\ Matthew}, Homily 50, \S\S\,4--5, trans.\ George Prevost and M.\ B.\ Riddle, in \textit{Nicene and Post-Nicene Fathers}, First Series, vol.\ 10 (CCEL dated electronic text, 2026-09-17; file header cites New York: Christian Literature Publishing Co., 1886).}
\end{quote}
